% Mention all the datasets used in the project and describe it

% =========================================================================== %
% --- Dataset Creation ---

\subsection{Raw Datasets} \label{subsection:raw}

Table \ref{table:dataset_characteristics} describes some relevant statistics of
the combined datasets used to train the models in the proposed framework. The
dataset contains images which pertain to three categories: authentic images,
fully AI-generated images and AI-manipulated (inpainted) images, all collected
from multiple publicly available benchmark datasets.

\begin{table}[h!t]
  \caption{Dataset Characteristics}
  \begin{center}
    \begin{tabularx}{\linewidth}{rYYYY}
      \toprule
        & \textbf{Column Name}                          & \textbf{Description}           & \textbf{Train Set Count} & \textbf{Test Set Count}\\
      \midrule
      1 & CIFAKE \cite{bird2023cifake}                  & Authentic and Generated Images & 100,000                  & 20,000 \\
      2 & Unbiased Tiny GenImage \cite{zhu2023genimage} & Authentic and Generated Images & 8,662                    & 2,165 \\
      \midrule
      3 & AutoSplice \cite{jia2023autosplice}           & Authentic and Inpainted Images & 5,519                    & 376 \\
      3 & CocoGlide \cite{guillaro2023cocoglide}        & Authentic and Inpainted Images & 819                      & 205 \\
      3 & SAGI \cite{giakoumoglou2025sagi}              & Authentic and Inpainted Images & 224,484                  & 12,688 \\
      \bottomrule
    \end{tabularx}
    \label{table:dataset_characteristics}
  \end{center}
\end{table}


\subsubsection{CIFAKE}
The CIFAKE dataset \cite{bird2023cifake} constitutes two classes:
\begin{itemize}
  \item{REAL - the authentic photographs}
  \item{FAKE - AI-generated synthetic images}
\end{itemize}
The primary use of this dataset was to improve the mode's ability to
differentiate between images which are real from images which are entirely
synthetic. CIFAKE produces strong images from  global image synthesis artifacts
produced by GAN and diffusion-based generators.

\subsubsection{Unbiased Tiny GenImage}
From the GenImage dataset \cite{zhu2023genimage}, the following subsets were
used:
\begin{itemize}
  \item{Nature - authentic images}
  \item{Midjourney - AI-generated images}
  \item{Glide - AI-generated images}
\end{itemize}
The purpose of this dataset was primarily to improve cross-generator
generalization. Since the generated images were from different diffusion
architecture origins, the model was able to learn the spectral and spatial
inconsistencies from the images rather than memorizing artifacts from a single
model.

\subsubsection{AutoSplice}
The AutoSplice dataset \cite{jia2023autosplice} was source for:
\begin{itemize}
  \item{Authentic - real images}
  \item{Forged\_JPEG90 - AI-manipulated images}
\end{itemize}
The focus of this particular dataset was on localized image manipulation and
splicing operations that helped the model learn discontinuities introduced
during selective image edition and region replacement.

\subsubsection{CocoGlide}
The CocoGlide dataset \cite{guillaro2023cocoglide} covered:
\begin{itemize}
  \item{Real - authentic images}
  \item{Fake - inpainted images}
\end{itemize}
The AI-manipulated samples are generated using GLIDE's inpainting techniques,
which was essential to the framework in teaching the model subtle boundary
inconsistencies and abnormal texture transitions, created in the AI-assisted
editing process.

\subsubsection{SAGI}
The SAGI dataset \cite{giakoumoglou2025sagi} contains subsets such as:
\begin{itemize}
  \item{original - authentic images}
  \item{brushnet - inpainted images}
  \item{controlnet - inpainted images}
  \item{removeanything - inpainted images}
\end{itemize}
The RAISE split of the dataset was the only category used during training. SAGI
introduced multiple modern inpainting approaches that contributed to the model's
ability to generalize across different pipelines and editing strategies.

\subsection{Processed Data} \label{subsection:processed}

Once the datasets were acquired, preprocessing was performed before training the
framework. The preprocessing pipeline included:
\begin{enumerate}
  \item{Image resizing and normalization}
  \item{Extraction of spectral features using Discrete Cosine Transform (DCT)}
  \item{Generation of monocular depth maps using MiDaS}
  \item{Feature vector construction for hybrid fusion}
\end{enumerate}

The processed dataset was stored in separate train and test directories. Each
sample contains a resized image file (.jpg), along with its corresponding
extracted feature files (.npy).

Most datasets were split using an 80-20 train-test ratio. The training split was
further divided into 80\% training and 20\% validation during model training.

\begin{table}[h!t]
  \caption{Label Distribution in the Combined Dataset}
  \begin{center}
    \begin{tabularx}{\linewidth}{rYcYY}
      \toprule
        & \textbf{Label}& \textbf{Numeric Label} & \textbf{Train Set Count} & \textbf{Test Set Count}\\
      \midrule
      1 & Authentic     & 0                      & 122,475                  & 15,226 \\
      2 & AI Generated  & 1                      & 54,797                   & 11,797 \\
      3 & Inpainted     & 2                      & 38,083                   & 2,494 \\
      \bottomrule
    \end{tabularx}
    \label{table:label_stats}
  \end{center}
\end{table}


The files follow a consistent naming scheme:
\verb|<label>_<original_file_name>|, where 0 indicates authentic images, 1
stands for AI-generated images and 2 is the label for inpainted images as shown
in table \ref{table:label_stats}. This preprocessing strategy essentially
ensures efficiency in feature loading during model training and subsequently
reduces repeated computation overhead.
